%%=============================================================================
%% Inleiding
%%=============================================================================

\chapter{\IfLanguageName{dutch}{Inleiding}{Introduction}}%
\label{ch:inleiding}

Aandachtstekortstoornis met hyperactiviteit, beter bekend als \textit{Attention Deficit Hyperactivity Disorder} (ADHD), is een neurobiologische ontwikkelingsstoornis die gekenmerkt wordt door aandachtsproblemen, impulsiviteit en/of hyperactiviteit. Naast deze kernsymptomen ervaren veel personen met ADHD moeilijkheden met executieve functies zoals planning, organisatie, tijdsinschatting en taakvoltooiing. Binnen een academische context, waar zelfstandige studieplanning en langdurige concentratie essentieel zijn, kunnen deze moeilijkheden leiden tot verhoogde stress, uitstelgedrag en verminderde studieprestaties. En dit is niet gelimiteerd tot enkel mensen met ADHD maar ook mensen met gelijkaardige symptomen.

Digitale planningsapplicaties worden vaak gebruikt als hulpmiddel voor taakbeheer en tijdsplanning. Hoewel er een breed aanbod aan tools bestaat, richten deze zich doorgaans op een algemene gebruikersgroep. Hierdoor sluiten ze niet altijd aan bij de specifieke cognitieve en emotionele noden van personen met ADHD. Interfaces zijn vaak complex, vereisen veel eigen structuur en maken gebruik van notificaties die als prikkelend of stresserend kunnen worden ervaren. Dit wijst op een kloof tussen bestaande digitale oplossingen en de concrete noden van studenten met ADHD.

Deze bachelorproef situeert zich binnen het domein van softwareontwikkeling en user experience design, en onderzoekt hoe een applicatie op maat van studenten met ADHD kan bijdragen aan beter taakbeheer en tijdsplanning. De focus ligt op het ontwerpen, implementeren en evalueren van een prototype waarin onder meer automatische taakopsplitsing, een afleidingsarme gebruikersinterface en rustige, motiverende feedbackmechanismen centraal staan.

\begin{itemize}
  \item context, achtergrond
  \item afbakenen van het onderwerp
  \item verantwoording van het onderwerp, methodologie
  \item probleemstelling
  \item onderzoeksdoelstelling
  \item onderzoeksvraag
  \item \ldots
\end{itemize}

\section{\IfLanguageName{dutch}{Probleemstelling}{Problem Statement}}%
\label{sec:probleemstelling}

Hoewel er tal van planningsapps beschikbaar zijn, blijken deze voor veel studenten met ADHD te complex en onvoldoende afgestemd op hun specifieke noden. Ze ervaren problemen met het structureren van grote taken, het inschatten van benodigde tijd en het volhouden van motivatie. Bestaande tools bieden vaak uitgebreide functionaliteiten, maar houden weinig rekening met beperkte werkgeheugencapaciteit, gevoeligheid voor prikkels en neiging tot uitstelgedrag.

De concrete doelgroep van dit onderzoek bestaat uit studenten met een (formeel gediagnosticeerde of zelfgerapporteerde) ADHD of gelijkaardige symptomen binnen Gent die digitale hulpmiddelen gebruiken voor studie- of werkgerelateerde taken. Het steekproefkader omvat studenten uit het hoger onderwijs (hogescholen en universiteiten) in Gent die bereid zijn deel te nemen aan bevragingen en gebruikerstesten.

De meerwaarde van dit onderzoek situeert zich op verschillende niveaus:
\begin{itemize}
  \item Studenten met ADHD in Gent die baat hebben bij een beter afgestemde planningsapplicatie;
  \item Studiebegeleiders en ondersteuningsdiensten binnen onderwijsinstellingen;
  \item Ontwikkelaars die inclusieve en cognitief toegankelijke software willen ontwerpen.
\end{itemize}

\section{\IfLanguageName{dutch}{Onderzoeksvraag}{Research question}}%
\label{sec:onderzoeksvraag}

De centrale onderzoeksvraag van deze bachelorproef luidt als volgt:

\begin{quote}
Hoe kan een digitale planningsflow studenten met ADHD of ADHD-symptomen ondersteunen bij het omzetten van academische deadlines naar haalbare acties, met aandacht voor cognitieve belasting, taakinitatie en herplanning 
\end{quote}

Om deze vraag te beantwoorden, wordt ze opgesplitst in deelvragen binnen twee domeinen.
\subsection*{Deelvragen probleemdomein}
\begin{itemize}
  \item Welke specifieke problemen ervaren studenten met ADHD bij planning en taakbeheer in het hoger onderwijs?
  \item Welke kenmerken van ADHD (bv. executieve functiestoornissen, tijdsblindheid, emotionele dysregulatie) hebben de grootste impact op studieplanning?
  \item Welke factoren veroorzaken stress en uitstelgedrag bij het uitvoeren van studietaken?
\end{itemize}

\subsection*{Deelvragen oplossingsdomein}
\begin{itemize}
  \item Welke UX/UI-principes dragen bij aan cognitieve toegankelijkheid en verminderen cognitieve belasting bij neurodiverse gebruikers?
  \item Hoe ervaren studenten een bescheiden prototype  van zo'n planningflow?
  \item Welke minimale functionaliteiten zijn noodzakelijk in een prototype om effectief ondersteuning te bieden bij dagelijkse planning?
\end{itemize}
\section{\IfLanguageName{dutch}{Onderzoeksdoelstelling}{Research objective}}%
\label{sec:onderzoeksdoelstelling}

 doelstelling van deze bachelorproef is het ontwikkelen en evalueren van een proof-of-concept van een planningsapplicatie op maat van studenten met ADHD.

Het beoogde resultaat is:
\begin{itemize}
  \item Een onderbouwde analyse van de noden en knelpunten van de doelgroep;
  \item Een werkend prototype (front-end met eenvoudige back-end) dat de voorgestelde ontwerpprincipes implementeert;
  \item Een evaluatie op basis van gebruikerstesten met studenten uit de doelgroep;
  \item Concrete aanbevelingen voor verdere ontwikkeling en opschaling.
\end{itemize}

De criteria voor succes zijn onder meer:
\begin{itemize}
  \item Een positieve beoordeling van de gebruiksvriendelijkheid door de doelgroep;
  \item Een gerapporteerde daling in ervaren cognitieve belasting tijdens gebruik;
  \item Indicaties van verbeterd overzicht en verhoogde motivatie of taakvoltooiing.
\end{itemize}
\section{\IfLanguageName{dutch}{Opzet van deze bachelorproef}{Structure of this bachelor thesis}}%
\label{sec:opzet-bachelorproef}

% Het is gebruikelijk aan het einde van de inleiding een overzicht te
% geven van de opbouw van de rest van de tekst. Deze sectie bevat al een aanzet
% die je kan aanvullen/aanpassen in functie van je eigen tekst.

De rest van deze bachelorproef is als volgt opgebouwd:

In Hoofdstuk~\ref{ch:stand-van-zaken} wordt een overzicht gegeven van de
bestaande literatuur rond ADHD, executieve functies en digitale
planningsondersteuning. Dit hoofdstuk kadert het probleem wetenschappelijk
en formuleert de ontwerpprincipes die als basis dienden voor het prototype.

In Hoofdstuk~\ref{ch:methodologie} wordt de gehanteerde onderzoeksaanpak
toegelicht. Het onderzoek verliep in vijf fasen: een literatuurstudie,
een vergelijkende app-analyse, semigestructureerde interviews, de
ontwikkeling van een klikbaar prototype en een kwantitatieve survey.

In Hoofdstuk~\ref{ch:requirementsanalyse} worden de functionele en
niet-functionele vereisten van de applicatie beschreven, afgeleid uit de
interviewresultaten en de vergelijkende app-analyse.

In Hoofdstuk~\ref{ch:ontwerp} wordt het ontwerpproces toegelicht: de
keuze voor de visuele stijl, de wireframes en de kernprincipes die het
ontwerp stuurden.

In Hoofdstuk~\ref{ch:implementatie} wordt de technische uitwerking van
het prototype beschreven. De applicatie werd gerealiseerd als een
standalone HTML-bestand en doorliep drie iteraties op basis van
opeenvolgende onderzoeksresultaten.

In Hoofdstuk~\ref{ch:functionaliteiten} worden de vijf schermen van het
prototype toegelicht: de dagweergave, het taakdetailscherm, het
invoerscherm, de herplanningsflow bij gemiste deadlines en de
patroonherkenningsmodule.

In Hoofdstuk~\ref{ch:surveyanalyse} worden de resultaten van de kwantitatieve
survey gerapporteerd en geïnterpreteerd. De survey werd ingevuld door 32
respondenten en valideerde zowel de probleemstelling als de gemaakte
ontwerpkeuzes.

In Hoofdstuk~\ref{ch:conclusie}, tenslotte, wordt de conclusie gegeven
en een antwoord geformuleerd op de centrale onderzoeksvraag. Daarbij
worden ook aanbevelingen geformuleerd voor verdere ontwikkeling en
toekomstig onderzoek binnen dit domein.